\documentclass[11pt]{article}
\usepackage{amsmath,amssymb,amsthm,hyperref}
\newtheorem{proposition}{Proposition}
\newcommand{\E}{\mathbb E}
\newcommand{\Pp}{\mathcal P}
\title{The nuclear error charge and a full-sampler obstruction to its mass relaxation}
\author{Campaign 3395: Engineer continuation}
\date{November 4, 2026}
\begin{document}
\maketitle

\paragraph{Status.}
T5 and the unrestricted fixed-factor lower bound remain unproved.
The new results below are an equivalent nuclear-error formulation of
T5 and an exact all-size counterexample to charging the
\emph{capped relative entrywise error} directly to $b_Q$.
The counterexample uses the full original sampler, one actual cut
that is entirely heavy, and thick conditional means with probability
tending to one. Its actual nuclear energy is exponentially small,
so it is not a counterexample to T5 or E4.

This is not another enlargement of the shrinking-accuracy window.
The strongest achieved approximation theorem remains reviewed V14.
Reviewer support for V1--V16, T1--T14, E1--E20, J8--J12, K1--K16,
the original cut, sampler and Boolean lift is reused without replay.
The nuclear/operator duality and PSD characterization in
arXiv:0706.4138v1, Proposition 2.1 and equation (2.6), retain their
independently confirmed attribution. No new external result is used.
The elementary symmetric recurrence used below is prior art already
present in the campaign's N14 and optional-support calculations;
the new ingredient is its use with fixed nonroot arcs retained by
the original random matching.

\section{An equivalent charge that keeps the nuclear norm}
Keep precisely the J2--J3 heavy matrices and V4's completions:
\[
 M_j=M_j(Q_j),\quad W_j^{\rm mat}=M_j(W_j),\quad
 L_j=M_j(H_j),\quad S_j^{\rm mat}=\min(R_j,3),\quad
 E'_j=S_j^{\rm mat}-M_j .
\]
Here $Q_j=A_jV_j$ is an actual cut product; the notation
$S_j^{\rm mat}$ is unrelated to its compulsory group set.
All these matrices have dimension $m_j=2^{k_j}$ and
$a_j=m_j^{-2}\sum W_j^{\rm mat}$.
Put
\[
 {\cal H}_j=\{k_j\ge5k_0/4\},\quad
 {\cal T}_j=\{a_j\ge2^{-k_0}\},\qquad k_0=N/2^{134},
\]
and define
\[
 \begin{split}
 {\mathfrak N}(Q)
 &=\sum_j\E\left[
     1_{{\cal H}_j\cap{\cal T}_j}
          \frac{\|E'_j\|_*^2}{m_j^3a_j}\right],\\
 {\mathfrak L}(Q)
 &=\sum_j\E\left[
     1_{{\cal H}_j\cap{\cal T}_j}
          \frac{\sum E'_j}{m_j^2a_j}\right].
 \end{split}                                                   \tag{Z1}
\]
The expectations use the unchanged original law and the original
unhonored-bit fixing for each actual cut. The indicators are
measurable in that conditioning. The positive threshold excludes
division by zero; use zero outside it.

\begin{proposition}
For admissible $N\ge2^{142}$ and every allowed circuit,
\[
 \begin{split}
 {\cal E}_{\rm thick}
 &\le3{\mathfrak N}(Q)+2^{16}s e^{-5k_0/32},\\
 {\mathfrak N}(Q)
 &\le3{\cal E}_{\rm thick}+2^{16}s e^{-5k_0/32}.
 \end{split}                                                   \tag{Z2}
\]
Thus T5 is equivalent, up to fixed constants, to
\[
 {\mathfrak N}(Q)\le Cb_Q+K M_n^\alpha(s+1)^\beta e^{-cN},
 \qquad C\ge0,\quad K,c,\beta>0,\quad\alpha\ge0.            \tag{Z3}
\]
Neither T5 nor (Z3) is proved here.
\end{proposition}
\begin{proof}
Use the two exact identities
$W_j^{\rm mat}=S_j^{\rm mat}-E'_j-L_j$ and
$E'_j=S_j^{\rm mat}-W_j^{\rm mat}-L_j$.
Squaring the nuclear triangle inequality bounds either energy by
three times the other plus
$27((k_j+1)^2+D_h)/(m_ja_j)$.
On ${\cal H}_j\cap{\cal T}_j$, reviewed V10 bounds the sum of these
remainders by $55323s e^{-5k_0/32}$.
For the first inequality only, add V8's strip
$k_0\le k_j<5k_0/4$, whose heavy energy is at most
$6s e^{-5k_0/32}$. Both coefficients are below $2^{16}$.
This proves (Z2).

In either direction an estimate of the displayed form transfers
with coefficient $3C$ and remainder constants
$3K+2^{16}$, $\alpha$, $\max(\beta,1)$, and
$\min(c,5/2^{139})$.
No new Cauchy--Schwarz budget or substituted law is used.
The reviewed common-law square-root transfer can then be applied
if (Z3) is eventually proved; stronger linear J10 is unnecessary.
\end{proof}

Because $0\le E'_j\le3$,
\[
 \|E'_j\|_*^2\le m_j\|E'_j\|_{\rm F}^2
                  \le3m_j\sum E'_j,\qquad
                         {\mathfrak N}\le3{\mathfrak L}. \tag{Z4}
\]
Therefore the following direct replacement for the V11--V12 loss
would be sufficient:
\[
 {\mathfrak L}(Q)\le Cb_Q+
                  K M_n^\alpha(s+1)^\beta e^{-cN}.       \tag{Z5}
\]
Unlike (Z3), assertion (Z5) is false. This remains so after capping
the errors, retaining only thick heavy matrices, and keeping all
original outer averaging.

\section{Fixed nonroot arcs seen by the original sampler}
Take any sufficiently large admissible $N\ge2^{154}$.
The base graph has maximum degree eight and at least $N-1$ edges.
For an edge $e$ let $Q_e$ be the set of edges incident to an
endpoint of $e$. At most $2\cdot8^2=128$ edges $f$ have
$Q_f\cap Q_e\ne\varnothing$, by the same neighborhood count used
in the original matching proof.
A greedy packing therefore permits a fixed set of
\[
                          v=\lfloor N/512\rfloor         \tag{Z6}
\]
edges with pairwise disjoint $Q_e$.
Orient each of these edges once and let $Y$ contain its
source-to-target variable $y_u$. Every $y_u$ has a nonroot target,
and their source groups are distinct.
Choose a fixed balanced partition $S,B$ of the $g=N+2$ groups
with all endpoints of these edges in $B$. This is possible since
$2v<g/2$. All choices precede the sampler.

Let $\ell=\ell(\eta)$ be the number of variables of $Y$ available
in the original common box $U_\eta$, and let $K=K(\omega)$ be the
number selected in the complete one-hot assignment. Thus
$0\le K\le\ell\le v$.

\begin{proposition}
Under the full original sampler, including the matching and its
prescribed order, with
\[
 \tau=N/2^{66},\qquad
 \delta_N=e^{-N/2^{19}}+2e^{-N/2^{68}},
\]
one has
\[
                 \Pr(K\ge\tau,\ \ell-K\ge\tau)
                              \ge1-\delta_N.             \tag{Z7}
\]
\end{proposition}
\begin{proof}
Retention events for the packed edges are independent because
their edge-mark neighborhoods are disjoint.
Each succeeds with probability at least $1/64$.
Since $v\ge N/1024$, their mean is at least $N/2^{16}$.
The reviewed elementary Bernoulli half-mean estimate shows that
at least $N/2^{17}$ of these edges are retained except with
probability $e^{-N/2^{19}}$.

The matching-preserving order lists endpoints of every retained
edge consecutively. At most $N/L$ such edges cross full-block
boundaries and at most $L$ are lost in the remainder.
Since $L=2^{23}\le N/2^{24}$, at least $N/2^{18}$ packed retained
edges lie inside full blocks.

Condition on any such matching and order. Let $t_i$ count these
oriented edges in block $i$. The event that the ordered slots
$(a_{i,1},b_{i,1})$ equal one of them has exact probability
$t_i/[L(L-1)]$. The other two slots can be any distinct remaining
vertices. If additionally $x_i=1$, the corresponding $y_u$ is
selected. If $x_i=0$, it is absent from the sample but present
in $U_\eta$: that source points to the next block rather than
its selected $b$ vertex.

For either choice of $x_i$, the indicators are independent across
blocks and have total mean at least
\[
 \sum_i\frac{t_i}{2L(L-1)}
                  \ge\frac{N}{2^{19}L^2}=N/2^{65}.
\]
Each success uses a distinct source. Applying the same half-mean
bound separately to the two signs gives counts at least
$N/2^{66}$, with total failure at most $2e^{-N/2^{68}}$.
Other appearances of $Y$ can only increase $K$ or $\ell-K$.
Averaging the actual matching and order proves (Z7).
No independently chosen active-source law has replaced the sampler.
\end{proof}

\section{A polynomial-size, entirely heavy actual cut}
Set
\[
 h=\lfloor2^{-150}N\rfloor,\quad d=h+1,\quad
 r=\lfloor N/2^{67}\rfloor,\quad
 \zeta=2^{-\lceil k_0/4\rceil},\quad
 \lambda=\frac{\binom rd}{\binom vd},\quad a_*=\zeta\lambda .
                                                               \tag{Z8}
\]
The stated range ensures $2d\le r<v$.
For each marked source group write its group sum as $y_u+z_u$,
where $z_u$ is the sum of all its other variables.
Let $U_k$ be the formal full-$B$ polynomial selecting exactly
$k$ variables from $Y$:
\[
 U_k=F_{B\setminus\operatorname{src}(Y)}
                         [t^k]\prod_{u\in\operatorname{src}(Y)}
                                      (z_u+t y_u).
\]
The coefficient notation is implemented by the ordinary positive
counting recurrence. Define
\[
 V=\sum_{k=r}^{v}\frac{U_k}{\binom kd},\qquad
 A=a_*F_S e_d(Y),\qquad Q=AV,                            \tag{Z9}
\]
where $e_d$ is the elementary symmetric polynomial of degree $d$.

\begin{proposition}
There is an original shared DAG of size $O(M_n+N^2)$ for (Z9),
with exactly one nonzero actual cut product, the displayed $AV$.
It is entirely its J2 heavy part, and
\[
              0\le Q\le\zeta F,\qquad b_Q\le\zeta\longrightarrow0
                        \quad\hbox{on the full Boolean cube}.  \tag{Z10}
\]
Its exact projection and sampled values are
\[
 \Pp Q=a_*F_S\sum_{k=r}^v U_k,\qquad
                         Q(X_\nu)=a_*1_{\{K\ge r\}}.     \tag{Z11}
\]
\end{proposition}
\begin{proof}
Compute all $U_k$ and $e_d(Y)$ by their positive dynamic programs,
using $O(M_n+N^2)$ gates. Constants such as $\binom kd^{-1}$
are allowed nonnegative real constants, here rational.
Build $V$ before $A$ and then their final product.
Since $d<v$, the intersection of supports of $e_d(Y)$ is empty,
so $G(A)=S$, while $G(V)=B$.
The last balanced gate is $A$; cutting it annihilates the output
and its exact divided context is $V$.
All further cut products vanish. Shared inputs do not change this.

Every monomial of $A$ has exactly $d>h$ optional nonroot source
groups outside $S$, and $V$ has none inside $S$. The entire one
actual cut, not an enumerated list of components, is therefore heavy.

On the Boolean cube, $e_d(Y)\le\binom vd$ and
$V\le F_B/\binom rd$, proving (Z10).
Full-group coverage handles all group-zero faces.
For a monomial of $U_k$, exactly $\binom kd$ multipliers in
$e_d(Y)$ are compatible, retaining every coefficient.
The reciprocal coefficient in $V$ cancels this count, proving
(Z11). The one-hot formula follows immediately.
\end{proof}

The amplitude is well above the thick threshold:
\[
 2^{-60d}\le\lambda\le2^{-57d},\qquad
             \lambda\ge2^{-k_0/512},\qquad
                              a_*\ge2^{-k_0/2}.          \tag{Z12}
\]
Indeed
$\lambda\ge((r-d+1)/v)^d\ge(r/(2v))^d\ge2^{-60d}$ and
$\lambda\le(r/v)^d\le2^{-57d}$.
Use $d\le N/2^{149}$ and $60/2^{15}<1/512$ for the next bound.
Finally $\lceil k_0/4\rceil\le k_0/4+1$ and $k_0\ge2^{20}$
give the last assertion.

\section{The mass relaxation diverges on thick histories}
For each original $(\eta,z)$ conditioning, let
$p=\Pr(K\ge r\mid\eta,z)$.
Every honored row source lies in $S$, whereas $Y$ and $V$ use
only $B$. Thus the full heavy matrix has identical rows, and
\[
                    a=a_*p,\qquad
 M(x,y)=a_*1_{\{K(y)\ge r\}}.                            \tag{Z13}
\]
The actual common-box completion has, on each column,
\[
 R(x,y)=a_*\frac{\binom\ell d}{\binom{K(y)}d}
                           1_{\{K(y)\ge r\}}\le\zeta .
                                                               \tag{Z14}
\]
Because $R\le\zeta\le1$, clipping at three does nothing:
$E'=R-M=E$. This is a bounded error of the very completion
used in V11--V12, not an alternative completion.

The conditional relative mass in (Z1) is consequently
\[
 \frac{\sum E'}{m^2a}
 =\frac{\E_{\rm honored}
    [1_{\{K\ge r\}}(\binom\ell d/\binom Kd-1)]}{p}.
                                                               \tag{Z15}
\]
On the event $\mathcal Z=\{K\ge\tau,\ell-K\ge\tau\}$ from (Z7),
one has $K\ge r$ and
\[
 \frac{\binom\ell d}{\binom Kd}
  =\prod_{i=0}^{d-1}\frac{\ell-i}{K-i}
  \ge(\ell/K)^d
  \ge(1+2^{-57})^d ,
                                                               \tag{Z16}
\]
since $K\le v\le N/512$.
The denominator $p\le1$ in (Z15) only strengthens this lower bound
after multiplying by the conditional probability of $\mathcal Z$.

Moreover $\Pr(p<1/2)\le2\Pr(K<r)\le2\delta_N$, by total
expectation in the original law. When $p\ge1/2$, (Z12) implies
$a\ge2^{-k_0}$.
The probability of $\neg{\cal H}$ is bounded by reviewed V8.
Thus, for $N\ge2^{154}$,
\[
 \Pr(\mathcal Z\cap{\cal H}\cap{\cal T})
 \ge1-3\delta_N-e^{-N/2^{22}}-e^{-5k_0/32}>1/2 .
                                                               \tag{Z17}
\]
For example each exponential has exponent at least eight, so
the total loss is below $11e^{-8}<1/2$.
All indicators here refer to the same full history;
none is averaged under a separately optimized distribution.

Equations (Z15)--(Z17) prove
\[
 {\mathfrak L}(Q)
 \ge\tfrac12((1+2^{-57})^d-1)
 \ge\tfrac12(e^{N/2^{208}}-1)\longrightarrow\infty .
                                                               \tag{Z18}
\]
For the second inequality use
$\log(1+t)\ge t/2$ on $0\le t\le1$ and $d>N/2^{150}$.
Since $M_n=\Theta(N)$, $s=O(M_n+N^2)$ and $b_Q\le\zeta\to0$,
the right side of any proposed (Z5) tends to zero.
This refutes (Z5) for every fixed allowed set of constants.

\section{Why the actual nuclear route survives}
The heavy matrix in (Z13) has rank at most one. For $a>0$ its
energy is exactly $a_*/m$, so
\[
             {\cal E}_{\rm thick}(Q)\le a_*2^{-k_0}.      \tag{Z19}
\]
The error in (Z14) also has identical rows.
On ${\cal H}\cap{\cal T}$ its nuclear norm equals its
Frobenius norm, whence
\[
 \frac{\|E'\|_*^2}{m^3a}
 \le\frac{\zeta^2}{ma}
 \le\zeta^2\,2^{-k_0/4}
 \le2^{-3k_0/4}.
 \quad\text{Thus}\quad
                       {\mathfrak N}(Q)\le2^{-3k_0/4}.   \tag{Z20}
\]
These are exponentially small under the full original sampler.
No nuclear norm has been replaced by entrywise mass in these
conclusions.

The old K11 obstruction had unbounded raw completion error and
zero heavy residual. Here the error is bounded by a number tending
to zero, the entire actual cut is heavy, and the thick event has
probability tending to one. This is a different obstruction to the
specific \emph{relative} mass relaxation after clipping.
The example is not an approximator of $R_n$; it tests the universal
$Q\le3F$ premise in T5 and the proposed sufficient shortcut.

Only the route through (Z5) is excluded. It would be incorrect to
retire T5, E4, (Z3), an adaptive amplitude-sensitive bound, or a
nuclear high-frequency argument from this construction.
In particular (Z2)--(Z3) retain the unrestricted nuclear structure
and the original-size remainder. Their missing defect charge is
still a proof obligation, not a theorem.

\paragraph{New exact replay.}
\path{relative_error_check.py} verifies a new rational circuit and
its sole entirely heavy cut, the coefficientwise projection,
all Boolean inputs on the four marked groups, and the actual
conditional matrices and completions.
An exact eight-square error Gram calculation verifies the rank-one
nuclear identities without interpreting entrywise mass as nuclear
cost. A separate bounded enumeration of the original ordered
four-tuple kernel at block size eight verifies the new matching-arc
trigger and its two bit probabilities.
The finite host is abstract, not an admissible physical graph or
an asymptotic original-sampler experiment.
All-size claims (Z7)--(Z20) rest on the proofs above, not on finite
extrapolation. New claims await independent review; no verified
novelty or publication qualification is asserted.
\end{document}
